% !TEX program = xelatex
% ============================================================
%  محضر الاجتماع — Meeting Minutes Template
%  نموذج لتوثيق محاضر اجتماعات المشروع بشكل منظم
%  Compile with: xelatex main.tex (run twice)
% ============================================================
%  Features:
%    - Header with date, time, location, chair, attendees, absentees
%    - Agenda items table
%    - Discussion summary section
%    - Decisions section
%    - Action items table (Action | Owner | Deadline | Status)
%    - Next meeting info
% ============================================================

\documentclass[11pt,a4paper]{article}

\usepackage[a4paper,margin=2.5cm,top=2.5cm,bottom=2.5cm,headheight=16pt]{geometry}
\usepackage{graphicx}
\usepackage{array}
\usepackage{booktabs}
\usepackage{tabularx}
\usepackage{multirow}
\usepackage{enumitem}
\usepackage{float}
\usepackage{titlesec}
\usepackage{fancyhdr}
\usepackage{setspace}
\usepackage{xcolor}
\usepackage{amsmath}

\usepackage{bidi}
\usepackage[no-math]{fontspec}
\usepackage[quiet]{polyglossia}

\setmainfont[Scale=1]{Amiri-Regular.ttf}
\newfontfamily\arabicfont[Script=Arabic,Scale=0.95]{Amiri-Regular.ttf}
\newfontfamily\englishfont[Scale=0.95]{TeX Gyre Termes}

\setdefaultlanguage[calendar=gregorian,hijricorrection=1]{arabic}
\setotherlanguage[variant=british]{english}

\usepackage[breaklinks,hidelinks]{hyperref}

\addto\captionsarabic{%
  \renewcommand{\contentsname}{الفهرس}%
  \renewcommand{\figurename}{الشكل}%
  \renewcommand{\tablename}{الجدول}%
}

\titleformat{\section}{\Large\bfseries}{\thesection}{0.7em}{}
\titlespacing*{\section}{0pt}{1.4ex plus 0.4ex}{0.9ex plus 0.3ex}

\pagestyle{fancy}
\fancyhf{}
\fancyhead[R]{\small\itshape محضر الاجتماع}
\fancyhead[L]{\small اسم المشروع}
\fancyfoot[C]{\small\thepage}
\renewcommand{\headrulewidth}{0.3pt}

\setlist[itemize]{topsep=0.3em, itemsep=0.15em, parsep=0pt}
\setlist[enumerate]{topsep=0.3em, itemsep=0.15em, parsep=0pt}

\newcolumntype{L}[1]{>{\raggedright\arraybackslash}p{#1}}
\newcolumntype{C}[1]{>{\centering\arraybackslash}p{#1}}

\setstretch{1.3}

\begin{document}

% ============================================================
% TITLE BLOCK
% ============================================================
\begin{center}
{\LARGE\bfseries محضر الاجتماع}\\[0.3em]
{\large Meeting Minutes}\\[1em]
\end{center}

% ============================================================
% MEETING INFO TABLE
% ============================================================
% TODO: املأ بيانات الاجتماع أدناه.
\begin{table}[H]
\centering
\renewcommand{\arraystretch}{1.4}
\begin{tabularx}{\textwidth}{L{3.5cm} X L{3.5cm} X}
\toprule
\textbf{التاريخ:} & --- & \textbf{الوقت:} & --- \\
\textbf{المكان:} & --- & \textbf{نوع الاجتماع:} & --- \\
\textbf{رئيس الاجتماع:} & --- & \textbf{محضر الاجتماع:} & --- \\
\textbf{اسم المشروع:} & --- & \textbf{رقم الاجتماع:} & --- \\
\bottomrule
\end{tabularx}
\end{table}

\vspace{0.5em}

% ============================================================
% ATTENDEES & ABSENTEES
% ============================================================
\begin{table}[H]
\centering
\renewcommand{\arraystretch}{1.3}
\begin{tabularx}{\textwidth}{L{3.5cm} X L{3.5cm} X}
\toprule
\textbf{الحضور:} & \multicolumn{3}{l}{%
\begin{minipage}[t]{\linewidth}
\vspace{0.2em}
\begin{itemize}
    \item الاسم الأول --- المنصب / القسم
    \item الاسم الثاني --- المنصب / القسم
    \item الاسم الثالث --- المنصب / القسم
    % TODO: أضف الحاضرين.
\end{itemize}
\vspace{-0.3em}
\end{minipage}} \\
\textbf{الغياب:} & \multicolumn{3}{l}{%
\begin{minipage}[t]{\linewidth}
\vspace{0.2em}
\begin{itemize}
    \item الاسم الرابع --- سبب الغياب
    % TODO: أضف الغائبين.
\end{itemize}
\vspace{-0.3em}
\end{minipage}} \\
\bottomrule
\end{tabularx}
\end{table}

% ============================================================
% 1. AGENDA ITEMS
% ============================================================
\section{بنود جدول الأعمال}
\label{sec:agenda}

% TODO: عدّل بنود جدول الأعمال حسب اجتماعك.
\begin{table}[H]
\centering
\caption{بنود جدول الأعمال}
\renewcommand{\arraystretch}{1.3}
\begin{tabularx}{\textwidth}{C{1cm} L{6cm} L{4cm} X}
\toprule
\textbf{رقم} & \textbf{البند} & \textbf{المسؤول} & \textbf{المدة} \\
\midrule
1 & مراجعة محضر الاجتماع السابق & رئيس الاجتماع & 10 دقائق \\
2 & عرض تقدم العمل & مدير المشروع & 20 دقيقة \\
3 & مناقشة المخاطر والمشكلات & فريق المشروع & 30 دقيقة \\
4 & مراجعة الميزانية & المسؤول المالي & 15 دقيقة \\
5 & القرارات المطلوبة & رئيس الاجتماع & 15 دقيقة \\
6 & جدولة الاجتماع القادم & الجميع & 5 دقائق \\
\bottomrule
\end{tabularx}
\end{table}

% ============================================================
% 2. DISCUSSION
% ============================================================
\section{المناقشات}
\label{sec:discussion}

% TODO: لخّص النقاط الرئيسية لكل بند من بنود جدول الأعمال.
\subsection{البند الأول: مراجعة محضر الاجتماع السابق}
تمت مراجعة محضر الاجتماع السابق والموافقة عليه مع تعديل بسيط. تم التأكد من إنجاز جميع بنود العمل السابقة.

\subsection{البند الثاني: عرض تقدم العمل}
تم عرض التقدم المحرز منذ الاجتماع السابق. المشروع يسير وفق الجدول المخطط له في معظم المهام، مع تأخير بسيط في مهمة واحدة.

\subsection{البند الثالث: مناقشة المخاطر والمشكلات}
تمت مناقشة المخاطر الرئيسية وطرق التعامل معها. تم تحديد خطر جديد يتعلق بتأخر المورد في التسليم.

% ============================================================
% 3. DECISIONS
% ============================================================
\section{القرارات}
\label{sec:decisions}

% TODO: اكتب القرارات المتخذة في الاجتماع.
\begin{enumerate}
    \item الموافقة على إضافة مورد بديل لتقليل خطر التأخير.
    \item تأجيل المهمة المتأخرة لمدة أسبوع مع إعادة توزيع الموارد.
    \item الموافقة على طلب ميزانية إضافية لشراء معدات جديدة.
\end{enumerate}

% ============================================================
% 4. ACTION ITEMS
% ============================================================
\section{بنود العمل}
\label{sec:actions}

% TODO: حدّد بنود العمل الناتجة عن الاجتماع.
\begin{table}[H]
\centering
\caption{بنود العمل الناتجة عن الاجتماع}
\renewcommand{\arraystretch}{1.3}
\begin{tabularx}{\textwidth}{C{1cm} L{5cm} L{3cm} L{2.5cm} L{2.5cm}}
\toprule
\textbf{رقم} & \textbf{بند العمل} & \textbf{المسؤول} & \textbf{الموعد النهائي} & \textbf{الحالة} \\
\midrule
1 & التواصل مع المورد البديل & مدير المشروع & --- & لم يبدأ \\
2 & إعادة جدولة المهمة المتأخرة & مدير المشروع & --- & قيد التنفيذ \\
3 & إعداد طلب الميزانية الإضافية & المسؤول المالي & --- & لم يبدأ \\
4 & تحديث سجل المخاطر & مدير المشروع & --- & لم يبدأ \\
5 & إرسال محضر الاجتماع للحضور & محضر الاجتماع & --- & لم يبدأ \\
\bottomrule
\end{tabularx}
\end{table}

% ============================================================
% 5. NEXT MEETING
% ============================================================
\section{الاجتماع القادم}
\label{sec:next}

% TODO: املأ تفاصيل الاجتماع القادم.
\begin{table}[H]
\centering
\renewcommand{\arraystretch}{1.4}
\begin{tabularx}{\textwidth}{L{3.5cm} X L{3.5cm} X}
\toprule
\textbf{التاريخ:} & --- & \textbf{الوقت:} & --- \\
\textbf{المكان:} & --- & \textbf{رئيس الاجتماع:} & --- \\
\bottomrule
\end{tabularx}
\end{table}

\vspace{1em}

% ============================================================
% SIGNATURES
% ============================================================
\noindent\rule{\textwidth}{0.4pt}

\vspace{0.5em}

\noindent
\begin{tabularx}{\textwidth}{X X}
\textbf{توقيع رئيس الاجتماع:} & \textbf{توقيع محضر الاجتماع:} \\[2em]
\rule{6cm}{0.3pt} & \rule{6cm}{0.3pt} \\
\end{tabularx}

\end{document}
